\documentclass[10pt,a4paper]{amsart}
\usepackage{amsmath,amssymb,amsthm,graphicx}
\usepackage{verbatim}
\newenvironment{reference}{\comment}{\endcomment}
%\newenvironment{reference}{}{}
\newenvironment{slogan}{\comment}{\endcomment}
\newenvironment{history}{\comment}{\endcomment}

% For commutative diagrams we use Xy-pic
\usepackage[all]{xy}

% We use 2cell for 2-commutative diagrams.
\xyoption{2cell}
% Theorem environments.
%
\theoremstyle{plain}
\newtheorem{theorem}[subsection]{Theorem}
\newtheorem{proposition}[subsection]{Proposition}
\newtheorem{lemma}[subsection]{Lemma}

\theoremstyle{definition}
\newtheorem{definition}[subsection]{Definition}
\newtheorem{example}[subsection]{Example}
\newtheorem{exercise}[subsection]{Exercise}
\newtheorem{situation}[subsection]{Situation}

\theoremstyle{remark}
\newtheorem{remark}[subsection]{Remark}
\newtheorem{remarks}[subsection]{Remarks}

\numberwithin{equation}{subsection}

% Macros
%
\def\lim{\mathop{\mathrm{lim}}\nolimits}
\def\colim{\mathop{\mathrm{colim}}\nolimits}
\def\Spec{\mathop{\mathrm{Spec}}}
\def\Hom{\mathop{\mathrm{Hom}}\nolimits}
\def\Ext{\mathop{\mathrm{Ext}}\nolimits}
\def\SheafHom{\mathop{\mathcal{H}\!\mathit{om}}\nolimits}
\def\SheafExt{\mathop{\mathcal{E}\!\mathit{xt}}\nolimits}
\def\Sch{\mathit{Sch}}
\def\Mor{\mathop{\mathrm{Mor}}\nolimits}
\def\Ob{\mathop{\mathrm{Ob}}\nolimits}
\def\Sh{\mathop{\mathit{Sh}}\nolimits}
\def\NL{\mathop{N\!L}\nolimits}
\def\CH{\mathop{\mathrm{CH}}\nolimits}
\def\proetale{{pro\text{-}\acute{e}tale}}
\def\etale{{\acute{e}tale}}
\def\QCoh{\mathit{QCoh}}
\def\Ker{\mathop{\mathrm{Ker}}}
\def\Im{\mathop{\mathrm{Im}}}
\def\Coker{\mathop{\mathrm{Coker}}}
\def\Coim{\mathop{\mathrm{Coim}}}

% Boxtimes
%
\DeclareMathSymbol{\boxtimes}{\mathbin}{AMSa}{"02}

%
% Macros for moduli stacks/spaces
%
\def\QCohstack{\mathcal{QC}\!\mathit{oh}}
\def\Cohstack{\mathcal{C}\!\mathit{oh}}
\def\Spacesstack{\mathcal{S}\!\mathit{paces}}
\def\Quotfunctor{\mathrm{Quot}}
\def\Hilbfunctor{\mathrm{Hilb}}
\def\Curvesstack{\mathcal{C}\!\mathit{urves}}
\def\Polarizedstack{\mathcal{P}\!\mathit{olarized}}
\def\Complexesstack{\mathcal{C}\!\mathit{omplexes}}
% \Pic is the operator that assigns to X its picard group, usage \Pic(X)
% \Picardstack_{X/B} denotes the Picard stack of X over B
% \Picardfunctor_{X/B} denotes the Picard functor of X over B
\def\Pic{\mathop{\mathrm{Pic}}\nolimits}
\def\Picardstack{\mathcal{P}\!\mathit{ic}}
\def\Picardfunctor{\mathrm{Pic}}
\def\Deformationcategory{\mathcal{D}\!\mathit{ef}}
\usepackage{cleveref}
\crefname{originalreference}{original reference}{original references}
\Crefname{originalreference}{Original reference}{Original references}
\makeatletter
\newlabel{example-rational-bridge}{{0E3M}{source}{}{sourcepointer-1}{}}
\newlabel{example-rational-bridge@cref}{{[originalreference][0][]0E3M}{[1][0][]source}{}{}{}}
\makeatother
\crefname{originalreference}{original reference}{original references}
\Crefname{originalreference}{Original reference}{Original references}
\makeatletter
\newlabel{lemma-closed-subscheme-nodal-curve}{{0E38}{source}{}{sourcepointer-2}{}}
\newlabel{lemma-closed-subscheme-nodal-curve@cref}{{[originalreference][0][]0E38}{[1][0][]source}{}{}{}}
\makeatother
\crefname{originalreference}{original reference}{original references}
\Crefname{originalreference}{Original reference}{Original references}
\makeatletter
\newlabel{lemma-closed-subscheme-reduced-gorenstein}{{0E34}{source}{}{sourcepointer-3}{}}
\newlabel{lemma-closed-subscheme-reduced-gorenstein@cref}{{[originalreference][0][]0E34}{[1][0][]source}{}{}{}}
\makeatother
\crefname{originalreference}{original reference}{original references}
\Crefname{originalreference}{Original reference}{Original references}
\makeatletter
\newlabel{lemma-facts-about-nodal-curves}{{0E37}{source}{}{sourcepointer-4}{}}
\newlabel{lemma-facts-about-nodal-curves@cref}{{[originalreference][0][]0E37}{[1][0][]source}{}{}{}}
\makeatother
\crefname{originalreference}{original reference}{original references}
\Crefname{originalreference}{Original reference}{Original references}
\makeatletter
\newlabel{lemma-rr}{{0BS6}{source}{}{sourcepointer-5}{}}
\newlabel{lemma-rr@cref}{{[originalreference][0][]0BS6}{[1][0][]source}{}{}{}}
\makeatother
\title{The canonical sheaf restricts to degree zero on a rational bridge}
\author{The Stacks Project Authors}
\date{stacks-project:curves.tex:local-snapshot-20261012; GNU FDL 1.2 or later}
\begin{document}
\maketitle
\noindent\textit{Editorial note.} All delivered mathematical text is the author's own, copied from the pinned Stacks Project chapter file sources/stacks-curves/pinned-original-source.tex; no statement or proof step was rewritten, reordered or removed.\par\medskip
\noindent\textit{Editorial note.} The counted proof is the author's complete argument: \par\medskip
\noindent\textit{Editorial note.} The label lemma-rational-bridge-zero names the counted statement, and the proof environment immediately following it in the source is delivered verbatim as its proof.\par\medskip
\noindent\textit{Editorial note.} The counted statement is carried by the author's own lemma environment in the chapter source.\par\medskip
\noindent\textit{Allowed standard prerequisites (editorial).}\begin{itemize}
\item The rational bridge setup. Same chapter, Example \textbackslash{}label\{example-rational-bridge\}. Proof lines 6321-6323 and 6327-6330: X' is the closed subscheme as in the example and k''/k'/k are the field extensions of the example; deg(omega\_C) = -2[k' : k] because C has genus 0, and deg(C \textbackslash{}cap X') = [k'' : k] = 2[k' : k].
\item Canonical sheaf on closed subschemes of nodal curves. Same chapter, Lemmas \textbackslash{}label\{lemma-closed-subscheme-reduced-gorenstein\}, \textbackslash{}label\{lemma-facts-about-nodal-curves\}, \textbackslash{}label\{lemma-closed-subscheme-nodal-curve\}. Proof lines 6323-6326: the short exact sequence 0 -> omega\_C -> omega\_X|\_C -> O\_\{C \textbackslash{}cap X'\} -> 0.
\item Riemann-Roch for curves. Same chapter, Lemma \textbackslash{}label\{lemma-rr\}. Proof lines 6327-6330: C has genus 0, so deg(omega\_C) = -2[k' : k].
\end{itemize}\medskip
\ifdefined\hypertarget\hypertarget{sourcepointer-1}{}\fi
\noindent\textit{Source reference.} Original reference 0E3M: sources/stacks-curves/pinned-original-source.tex, source lines 6201--6201. Prerequisite pointer invoked at proof line(s) 6321; it names a standard fact of the same chapter that the counted argument uses, and every step of the counted proof built on it is retained verbatim in the delivered body.\par\medskip
\ifdefined\hypertarget\hypertarget{sourcepointer-2}{}\fi
\noindent\textit{Source reference.} Original reference 0E38: sources/stacks-curves/pinned-original-source.tex, source lines 4794--4794. Prerequisite pointer invoked at proof line(s) 6332; it names a standard fact of the same chapter that the counted argument uses, and every step of the counted proof built on it is retained verbatim in the delivered body.\par\medskip
\ifdefined\hypertarget\hypertarget{sourcepointer-3}{}\fi
\noindent\textit{Source reference.} Original reference 0E34: sources/stacks-curves/pinned-original-source.tex, source lines 1012--1012. Prerequisite pointer invoked at proof line(s) 6330; it names a standard fact of the same chapter that the counted argument uses, and every step of the counted proof built on it is retained verbatim in the delivered body.\par\medskip
\ifdefined\hypertarget\hypertarget{sourcepointer-4}{}\fi
\noindent\textit{Source reference.} Original reference 0E37: sources/stacks-curves/pinned-original-source.tex, source lines 4763--4763. Prerequisite pointer invoked at proof line(s) 6331; it names a standard fact of the same chapter that the counted argument uses, and every step of the counted proof built on it is retained verbatim in the delivered body.\par\medskip
\ifdefined\hypertarget\hypertarget{sourcepointer-5}{}\fi
\noindent\textit{Source reference.} Original reference 0BS6: sources/stacks-curves/pinned-original-source.tex, source lines 1218--1218. Prerequisite pointer invoked at proof line(s) 6335; it names a standard fact of the same chapter that the counted argument uses, and every step of the counted proof built on it is retained verbatim in the delivered body.\par\medskip
\begin{lemma}
\label{lemma-rational-bridge-zero}
Let $k$ be a field. Let $X$ be a proper scheme over $k$ having dimension $1$
and $H^0(X, \mathcal{O}_X) = k$. Assume the singularities of $X$ are
at-worst-nodal. Let $C \subset X$ be a rational bridge
(Example \ref{example-rational-bridge}). Then $\deg(\omega_X|_C) = 0$.
\end{lemma}

\begin{proof}
Let $X' \subset X$ be as in the example. Then we have a short exact
sequence
$$
0 \to \omega_C \to \omega_X|_C \to \mathcal{O}_{C \cap X'} \to 0
$$
See Lemmas \ref{lemma-closed-subscheme-reduced-gorenstein},
\ref{lemma-facts-about-nodal-curves}, and
\ref{lemma-closed-subscheme-nodal-curve}.
With $k''/k'/k$ as in the example we see that
$\deg(\omega_C) = -2[k' : k]$ as $C$ has genus $0$
(Lemma \ref{lemma-rr})
and $\deg(C \cap X') = [k'' : k] = 2[k' : k]$.
Hence $\deg(\omega_X|_C) = 0$.
\end{proof}

\end{document}
